
\documentclass{article}
\usepackage{colortbl}
\usepackage{makecell}
\usepackage{multirow}
\usepackage{supertabular}

\begin{document}

\newcounter{utterance}

\centering \large Interaction Transcript for game `cladder', experiment `full\_v1.5\_default', episode 9436 with score{-}Qwen3.5{-}9B{-}sp.
\vspace{24pt}

{ \footnotesize  \setcounter{utterance}{1}
\setlength{\tabcolsep}{0pt}
\begin{supertabular}{c@{$\;$}|p{.15\linewidth}@{}p{.15\linewidth}p{.15\linewidth}p{.15\linewidth}p{.15\linewidth}p{.15\linewidth}}
    \# & \multicolumn{2}{c}{Player} && \multicolumn{2}{c}{Game Master} \\
    \hline

    \theutterance \stepcounter{utterance}  
    & & & \multicolumn{4}{p{0.6\linewidth}}{
        \cellcolor[rgb]{0.9,0.9,0.9}{
            \makecell[{{p{\linewidth}}}]{
                \texttt{\tiny{[P1$\langle$GM]}}
                \texttt{You are an expert in causal inference. The following question is not a typical commonsense query, but rather a meticulously designed question created by a professor specializing in causal inference, intended to assess the students' mastery of the course content.} \\
\texttt{Imagine a self{-}contained, hypothetical world with only the following conditions, and without any unmentioned factors or causal relationships: Rixq has a direct effect on zuph and xevu. Swoy has a direct effect on zuph and xevu. Zuph has a direct effect on xevu. Swoy is unobserved. We know that rixq or swoy causes zuph. rixq and swoy and zuph causes xevu. We observed an individual is not swoy. Would an individual is xevu if rixq instead of not rixq?} \\
\texttt{Start your answer with 'Yes' or 'No', followed by additional reasoning or evidence to support your explanation.} \\
            }
        }
    }
     \\ \\

    \theutterance \stepcounter{utterance}  
    & \multicolumn{4}{p{0.6\linewidth}}{
        \cellcolor[rgb]{0.561,0.737,0.561}{
            \makecell[{{p{\linewidth}}}]{
                \texttt{\tiny{[P1$\rangle$GM]}}
                \texttt{No} \\
\\ 
\texttt{**Reasoning:**} \\
\\ 
\texttt{To determine if an individual would be \$xevu\$ if \$rixq\$ were true instead of false, we must analyze the causal structure and the specific state of the individual in question.} \\
\\ 
\texttt{1.  **Causal Structure Analysis**:} \\
\texttt{    *   \$rixq \textbackslash{}rightarrow zuph\$} \\
\texttt{    *   \$swoy \textbackslash{}rightarrow zuph\$} \\
\texttt{    *   \$zuph \textbackslash{}rightarrow xevu\$} \\
\texttt{    *   \$rixq \textbackslash{}rightarrow xevu\$ (Direct effect)} \\
\texttt{    *   \$swoy \textbackslash{}rightarrow xevu\$ (Direct effect)} \\
\texttt{    *   \$swoy\$ is unobserved.} \\
\\ 
\texttt{2.  **The Individual's State**:} \\
\texttt{    *   We are told the observed individual is **not \$swoy\$** (\$swoy = 0\$).} \\
\texttt{    *   The question asks for the counterfactual: Would this individual be \$xevu\$ if \$rixq = 1\$ (instead of \$rixq = 0\$)?} \\
\\ 
\texttt{3.  **Evaluating the Outcome (\$xevu\$)**:} \\
\texttt{    *   The outcome \$xevu\$ is caused by \$zuph\$, \$rixq\$, and \$swoy\$.} \\
\texttt{    *   Since \$swoy = 0\$, the contribution of \$swoy\$ to \$xevu\$ is null} \\
            }
        }
    }
    & & \\ \\

    \theutterance \stepcounter{utterance}  
    & & & \multicolumn{2}{p{0.3\linewidth}}{
        \cellcolor[rgb]{0.9,0.9,0.9}{
            \makecell[{{p{\linewidth}}}]{
                \texttt{\tiny{[GM$|$GM]}}
                \texttt{No} \\
            }
        }
    }
    & & \\ \\

    \theutterance \stepcounter{utterance}  
    & & & \multicolumn{2}{p{0.3\linewidth}}{
        \cellcolor[rgb]{0.9,0.9,0.9}{
            \makecell[{{p{\linewidth}}}]{
                \texttt{\tiny{[GM$|$GM]}}
                \texttt{no} \\
            }
        }
    }
    & & \\ \\

    \theutterance \stepcounter{utterance}  
    & & & \multicolumn{2}{p{0.3\linewidth}}{
        \cellcolor[rgb]{0.9,0.9,0.9}{
            \makecell[{{p{\linewidth}}}]{
                \texttt{\tiny{[GM$|$GM]}}
                \texttt{game\_result = WIN} \\
            }
        }
    }
    & & \\ \\

\end{supertabular}
}

\end{document}
